\section{Theta lifts and one-mode coercivity}
\label{sec:theta-coercivity}

\subsection{The theta image and its polynomial hierarchy}

Use the Fourier convention $\widehat f(t)=\int_{\mathbb R}f(x)e^{-itx}\,dx$. Put

\[
 h(u)=(4\pi^2u^4-6\pi u^2)e^{-\pi u^2},\qquad
 (\Theta\psi)(x)=e^{x/2}\sum_{n\ge1}\psi(ne^x),\qquad K=\Theta h.
\]

The integral defining the Mellin transform of $h$ gives, exactly,

\begin{equation}
M_h(s):=\int_0^\infty h(u)u^{s-1}\,du
   =\frac12s(s-1)\pi^{-s/2}\Gamma(s/2).
 \label{eq:theta-hier-1}
\end{equation}

Indeed the two Gaussian moments are
$2\pi^{-s/2}\Gamma(s/2+2)$ and
$-3\pi^{-s/2}\Gamma(s/2+1)$; the Gamma recurrence gives \eqref{eq:theta-hier-1}.
For $\operatorname{Re}s>1$, the Mellin sum is absolutely convergent and gives
$\zeta(s)M_h(s)=\xi(s)$. Its Fourier restriction at
$s=1/2-it$, continued by the functional equation, is
$\widehat K(t)=\Xi(t)$. The exact Poisson relation for the self-Fourier function $h$, with $h(0)=\int h=0$, makes $K$ even.

Let $E=u\partial_u+1/2$. Termwise differentiation, justified by Gaussian decay at infinity, gives the exact intertwining

\begin{equation}
\partial_x\Theta\psi=\Theta E\psi.
 \label{eq:theta-hier-2}
\end{equation}

Thus $K^{(j)}=\Theta E^jh$ and
$\widehat{K^{(j)}}(t)=(it)^j\Xi(t)$. These profiles all have zero sampling vector under the actual map

\begin{equation}
\mathcal S f=(\sqrt{m_\gamma}\,\widehat f(\gamma))_{\gamma\in\Gamma}.
 \label{eq:theta-hier-3}
\end{equation}

Here the map is defined on rapidly decreasing smooth profiles for which the displayed sum converges. Its squared norm is the extended Weil seminorm. Compact cutoffs converge in that seminorm by the sampling bound proved below; consequently this map agrees with the completion of compact test functions. In particular, these nonzero physical profiles have image zero under \eqref{eq:theta-hier-3}. 

As $x\to+\infty$, with $u=e^x$, the first theta summand dominates, while the sum of all other terms and each fixed derivative is exponentially smaller. Hence

\begin{equation}
K^{(j)}(x)=4\pi^2(-2\pi)^j
       u^{9/2+2j}e^{-\pi u^2}(1+O_j(u^{-2})).
 \label{eq:theta-hier-4}
\end{equation}

The relative $u^{-2}$ error comes from the lower polynomial in $h$ and differentiation; terms $n\ge2$ have the additional exponential factor $e^{-3\pi u^2}$ times a fixed polynomial. Integration after the change of variable $u=e^x$ gives

\begin{equation}
\int_L^\infty|K^{(j)}(x)|^2\,dx
  \sim4\pi^3(2\pi)^{2j}
         \lambda^{7+4j}e^{-2\pi\lambda^2},\qquad\lambda=e^L.
 \label{eq:theta-hier-5}
\end{equation}

This follows from
$\int_\lambda^\infty u^a e^{-2\pi u^2}du
\sim\lambda^{a-1}e^{-2\pi\lambda^2}/(4\pi)$.

\subsection{Division at an actual zero: the exact first-order lift}

Fix an actual positive ordinate $\gamma$, of its actual multiplicity $m_\gamma$, and put $s_\gamma=1/2-i\gamma$. Define

\begin{equation}
(R_\gamma K)(x)
   =-i e^{i\gamma x}\int_x^\infty e^{-i\gamma y}K(y)\,dy.
 \label{eq:theta-hier-6}
\end{equation}

Differentiation gives
$(-i\partial_x-\gamma)R_\gamma K=K$.
The identity $\widehat K(\gamma)=0$ also gives the equivalent left-tail formula
$R_\gamma K=i e^{i\gamma x}\int_{-\infty}^x e^{-i\gamma y}K(y)dy$.
Both ends therefore decay faster than every exponential, and Fourier transformation yields

\begin{equation}
\widehat{R_\gamma K}(t)=\frac{\Xi(t)}{t-\gamma},
 \label{eq:theta-hier-7}
\end{equation}

with the entire extension at $\gamma$.

The corresponding raw theta input can be written directly:

\begin{equation}
\psi_\gamma(u)
  =-i u^{-s_\gamma}\int_u^\infty h(v)v^{s_\gamma-1}\,dv.
 \label{eq:theta-hier-8}
\end{equation}

It satisfies

\begin{equation}
(-iE-\gamma)\psi_\gamma=h,
 \qquad \Theta\psi_\gamma=R_\gamma K.
 \label{eq:theta-hier-9}
\end{equation}

The first identity follows by differentiation. The theta series and its derivative converge locally normally because the upper tail of \eqref{eq:theta-hier-8} is Gaussian. Equation \eqref{eq:theta-hier-2} shows that its theta image satisfies the same first-order equation as \eqref{eq:theta-hier-6}, and both decay at positive infinity. Their difference is a multiple of $e^{i\gamma x}$; positive-tail decay forces that multiple to be zero. This proves the second identity.

The two raw endpoint expansions are

\begin{equation}
\psi_\gamma(u)=-iM_h(s_\gamma)u^{-s_\gamma}+O_\gamma(u^2)
       \quad(u\downarrow0),
 \qquad
 \psi_\gamma(u)\sim\frac{-i}{2\pi u^2}h(u)
       \quad(u\to\infty).
 \label{eq:theta-hier-10}
\end{equation}

For the first, split the integral in \eqref{eq:theta-hier-8} at zero and use $h(v)=O(v^2)$. Its leading coefficient is nonzero by \eqref{eq:theta-hier-1}: $s_\gamma$ is neither $0$ nor $1$, and the Gamma function has no zeros. For the second, integrate the Gaussian tail once, or apply the elementary Laplace-tail asymptotic to each polynomial term. Thus the raw lift has a nonzero $u^{-1/2+i\gamma}$ term at zero and its squared modulus is not integrable there. The theta image nevertheless has the two rapidly decreasing tails proved directly by \eqref{eq:theta-hier-6}.

The exact Mellin relation behind this cancellation is

\begin{equation}
M_{\psi_\gamma}(s)=\frac{M_h(s)}{i(s-s_\gamma)},\qquad
 \zeta(s)M_{\psi_\gamma}(s)=\frac{\xi(s)}{i(s-s_\gamma)}.
 \label{eq:theta-hier-11}
\end{equation}

The first holds initially in the strip $\operatorname{Re}s>1/2$, where the endpoint integrations converge. The second holds initially for $\operatorname{Re}s>1$ by the absolutely convergent theta sum. The zero of $\zeta$ at $s_\gamma$ cancels the displayed raw Mellin pole. Equations \eqref{eq:theta-hier-6}--\eqref{eq:theta-hier-11} prove the relation directly, including the domain and the coefficient of the singular term.

\subsection{Repeated and paired divisors}

For $1\le r\le m_\gamma$, repeated application of \eqref{eq:theta-hier-6} gives

\begin{equation}
F_{\gamma,r}(x)=\frac{(-i)^r e^{i\gamma x}}{(r-1)!}
   \int_x^\infty (y-x)^{r-1}e^{-i\gamma y}K(y)\,dy,
 \qquad
 \widehat F_{\gamma,r}(t)=\frac{\Xi(t)}{(t-\gamma)^r}.
 \label{eq:theta-hier-12}
\end{equation}

The moment identities $\Xi^{(j)}(\gamma)=0$, $0\le j<r$, convert \eqref{eq:theta-hier-12} to the corresponding left-tail integral, proving decay at both ends. Its actual sampling norm is

\begin{equation}
q(F_{\gamma,r})=
 \begin{cases}
 0,&r<m_\gamma,\\
 m_\gamma\left|\Xi^{(m_\gamma)}(\gamma)/m_\gamma!\right|^2,
       &r=m_\gamma.
 \end{cases}
 \label{eq:theta-hier-13}
\end{equation}

This is obtained simply by evaluating \eqref{eq:theta-hier-12}'s entire transform at every signed ordinate: all values vanish except the final possible value at $\gamma$.

For the even real profiles define

\begin{equation}
\widehat G_{\gamma,r}(t)=\frac{\Xi(t)}{(t^2-\gamma^2)^r},
 \qquad 1\le r\le m_\gamma.
 \label{eq:theta-hier-14}
\end{equation}

The transform is real and even on the real axis, so its inverse is real and even. Applying $R_\gamma^rR_{-\gamma}^r$ gives this profile, with the same two-tail argument. At the two removed zeros the exact values are equal, so

\begin{equation}
q(G_{\gamma,r})=
 \begin{cases}
 0,&r<m_\gamma,\\
 2m_\gamma\left|
 \dfrac{\Xi^{(m_\gamma)}(\gamma)}{m_\gamma!(2\gamma)^{m_\gamma}}
                  \right|^2,&r=m_\gamma.
 \end{cases}
 \label{eq:theta-hier-15}
\end{equation}

For each application of $R_\beta$, the positive-tail Laplace integral multiplies its input's leading term by $-i/(2\pi u^2)$. Therefore

\begin{equation}
F_{\gamma,r}(x)\sim
   4\pi^2\left(\frac{-i}{2\pi}\right)^r
       u^{9/2-2r}e^{-\pi u^2},
 \qquad
 G_{\gamma,r}(x)\sim
   4\pi^2\frac{(-1)^r}{(2\pi)^{2r}}
       u^{9/2-4r}e^{-\pi u^2}.
 \label{eq:theta-hier-16}
\end{equation}

In particular $G_{\gamma,1}(x)\sim-e^{x/2}e^{-\pi e^{2x}}$, while $K(x)\sim4\pi^2e^{9x/2}e^{-\pi e^{2x}}$. The ratio is exactly asymptotic to $-e^{-4x}/(4\pi^2)$. Correspondingly,

\begin{equation}
\begin{aligned}
 \int_L^\infty |F_{\gamma,r}|^2
   &\sim4\pi^3(2\pi)^{-2r}\lambda^{7-4r}e^{-2\pi\lambda^2},\\
 \int_L^\infty |G_{\gamma,r}|^2
   &\sim4\pi^3(2\pi)^{-4r}\lambda^{7-8r}e^{-2\pi\lambda^2}.
\end{aligned}
 \label{eq:theta-hier-17}
\end{equation}

The first paired raw lift is explicitly
$\psi_{\gamma,\mathrm{pair}}=(\psi_\gamma-\psi_{-\gamma})/(2\gamma)$. Its Mellin transform is

\begin{equation}
\frac{M_h(s)}{-(s-1/2)^2-\gamma^2},
 \label{eq:theta-hier-18}
\end{equation}

and its theta image is $G_{\gamma,1}$. The nonzero endpoint terms at zero have exponents $-1/2\pm i\gamma$. Their squared integral diverges: after $u=e^{-v}$, the diagonal terms have a positive constant mean, while the cross term has a bounded oscillatory integral. For a repeated paired divisor, the raw Mellin poles have order $r$, giving nonzero polynomial factors in $\log u$ of degree $r-1$ at those two exponents. The same conclusion follows because the diagonal integral grows like $V^{2r-1}$, while integration by parts bounds the oscillatory cross term by $O(V^{2r-2})$. Thus the precise theta lifts carry explicitly quantified endpoint singularities; applying the square-integrable prolate estimates to them requires an additional truncation estimate for those terms.

\subsection{What direct compact truncation proves}

The local zero-count bound, including multiplicities, is
$\#\{\gamma:n\le\gamma<n+1\}\ll\log(2+|n|)$. Here is a direct justification of the bound used. For large positive $n$, apply Jensen's formula to $\zeta(s)$ on the disk of radius $3$ centered at $2+i(n+1/2)$; the disk of radius $5/2$ contains the whole zero strip with ordinates in $[n,n+1]$. The value at the center has modulus at least $1/\zeta(2)$, by the absolutely convergent reciprocal Dirichlet series. Euler--Maclaurin summation at an integer of size $n$, retaining the Bernoulli terms through order $4$, bounds the modulus throughout the outer disk by $O(n^2)$: on $-1\le\operatorname{Re}s\le5$, the finite sum is $O(n^2)$, and the remainder is bounded by a constant times $|s|^4 n^{-\operatorname{Re}s-3}=O(n^2)$. Jensen's formula gives the claimed logarithmic count. Conjugation treats negative ordinates, and the finitely many remaining intervals are absorbed in the constant. Unit-interval Sobolev evaluation therefore gives

\begin{equation}
q(f)\ll\int_{\mathbb R}\log(2+|t|)
      (|\widehat f(t)|^2+|\widehat f'(t)|^2)\,dt.
 \label{eq:theta-hier-19}
\end{equation}

For $\kappa\ge1$, use
$\log(2+|t|)\le\log(2+\kappa)+t^2/\kappa^2$. Parseval gives the concrete physical bound

\begin{equation}
q(f)\ll \log(2+\kappa)(\|f\|_2^2+\|xf\|_2^2)
     +\kappa^{-2}(\|f'\|_2^2+\|(xf)'\|_2^2).
 \label{eq:theta-hier-20}
\end{equation}

Choose a fixed smooth transition and an even cutoff $\chi_L$ supported in $[-L,L]$, equal to one on $[-L+\lambda^{-2},L-\lambda^{-2}]$, with $|\chi_L'|\ll\lambda^2$. Apply \eqref{eq:theta-hier-20}, with $\kappa=\lambda^2$, to $(1-\chi_L)f$. Equations \eqref{eq:theta-hier-4}, \eqref{eq:theta-hier-16}, and the corresponding differentiated estimates show that a profile with leading positive-tail exponent $u^a e^{-\pi u^2}$ obeys

\begin{equation}
q((1-\chi_L)f)\ll_f
      (\log\lambda)^3\lambda^{2a-2}e^{-2\pi\lambda^2}.
 \label{eq:theta-hier-21}
\end{equation}

The transition starts at $L-\lambda^{-2}$; replacing $\lambda$ by $\lambda e^{-\lambda^{-2}}$ changes the exponential by a bounded factor. The derivative of the cutoff has exactly the $\lambda^2$ size absorbed by $\kappa^{-2}$, and the factor $x$ contributes at most $O(L)$ in these tail integrals. This proves \eqref{eq:theta-hier-21}, including the transition interval.

Consequently,

\begin{equation}
q(\chi_L K^{(j)})\ll_j
      (\log\lambda)^3\lambda^{7+4j}e^{-2\pi\lambda^2}.
 \label{eq:theta-hier-22}
\end{equation}

For every profile \eqref{eq:theta-hier-12} or \eqref{eq:theta-hier-14}, the seminorm triangle inequality gives
$|\sqrt{q(\chi_L f)}-\sqrt{q(f)}|\le\sqrt{q((1-\chi_L)f)}$; substitute the exact values \eqref{eq:theta-hier-13}, \eqref{eq:theta-hier-15} and exponents from \eqref{eq:theta-hier-16}. In particular the zero entries in \eqref{eq:theta-hier-15} yield the upper scale
$(\log\lambda)^3\lambda^{7-8r}e^{-2\pi\lambda^2}$ for their compact approximants. These identities verify the tail hierarchy. The proved compact-cutoff bound has exponential $e^{-2\pi\lambda^2}$; the sharper original prolate construction has $e^{-4\pi\lambda^2}$. The exact maps above relate the profiles, but \eqref{eq:theta-hier-21} does not give a prolate-scale lower bound.

\subsection{Exact removal of the stated trial mode}

Let $v_L$ be the actual normalized even prolate trial, and define

\begin{equation}
b_L=\inf\{q(f):f\in\operatorname{Dom}q_L,
       \|f\|_2=1,\ \langle f,v_L\rangle_2=0\}.
 \label{eq:theta-hier-23}
\end{equation}

The exact comparison is

\begin{equation}
a_1(L)+a_2(L)-q(v_L)\le b_L\le a_2(L).
 \label{eq:theta-hier-24}
\end{equation}

For the lower inequality, every admissible $(v_L,f)$ is an orthonormal pair. Ky Fan's trace inequality gives $q(v_L)+q(f)\ge a_1+a_2$. For the upper inequality, the span of the first two eigenvectors meets $v_L^\perp$ nontrivially, and the Rayleigh quotient on that span is at most $a_2$. In particular $a_2-q(v_L)\le b_L\le a_2$, since $a_1\ge0$. This explicitly relates removal of the stated theta/prolate mode to the desired denominator; it needs no separate assertion that the trial aligns with the first eigenvector.

The remaining lower estimate concerns the actual restriction
\eqref{eq:theta-hier-23}:
\[
 b_L\ge c\lambda^{9+\alpha}e^{-4\pi\lambda^2},
 \qquad L=\log\lambda,
\]
with fixed $c>0$ and $\alpha>0$ along an unbounded sequence of
windows.  This estimate remains open.  Equation~\eqref{eq:theta-hier-24}
places $b_L$ within the already controlled trial energy of $a_2(L)$.
Equations~\eqref{eq:theta-hier-8} and \eqref{eq:theta-hier-18} identify
the raw endpoint terms that an argument using these theta lifts must
control at the displayed scale.
